\documentclass[11pt]{article}

\usepackage[margin=1in]{geometry}
\usepackage{amsmath, amssymb, amsfonts}
\usepackage{enumitem}
\usepackage{hyperref}

\title{Dimensional Lift and Advanced Filter Gates}
\author{kinematic-classifier-sandbox}
\date{\today}

\begin{document}

\maketitle

\section{Purpose}

This document explains the transition logic from 1D witness problems toward
higher-dimensional work and documents why IMM, PF, and RBPF remain
decision-gated rather than automatically implemented.

\section{Problem And Notation}

This note has two proof obligations:
\begin{equation}
    \text{Can the framework survive a lift from scalar 1D trajectories to vector-valued trajectories?}
\end{equation}
and
\begin{equation}
    \text{What measured failure evidence would justify IMM, PF, or RBPF?}
\end{equation}

The most important symbols are:
\begin{itemize}[leftmargin=1.5em]
    \item \(a_i\): one scalar-assumption audit row,
    \item \(z_t \in \mathbb{R}^d\): vector observation at time \(t\),
    \item \(\Delta\): measured gain between two ladder rungs,
    \item \(x_t = (u_t, v_t)\): sampled-plus-analytic latent split for RBPF reasoning.
\end{itemize}

\section{Dimensional Lift Audit}

The main dimensional-readiness surfaces are:
\begin{itemize}[leftmargin=1.5em]
    \item \texttt{dimensional\_lift\_audit.py}
    \item \texttt{pca\_dimensionality\_audit.py}
\end{itemize}

These modules classify code as:
\begin{itemize}[leftmargin=1.5em]
    \item dimension-agnostic,
    \item adapter-compatible,
    \item rewrite-required.
\end{itemize}

The operational meaning is:
\begin{equation}
    \text{module status}
    \in
    \{
        \text{agnostic},
        \text{adapter-compatible},
        \text{rewrite-required}
    \},
\end{equation}
where the label is determined by whether the implementation depends on scalar
measurement semantics, scalar state layout, or scalar-derived feature logic.

\section{Scalar-Assumption Inventory Logic}

The audit looks for assumptions such as:
\begin{itemize}[leftmargin=1.5em]
    \item scalar measurements,
    \item scalar state dimensions,
    \item scalar derivative construction,
    \item scalar-only covariance or plotting expectations.
\end{itemize}

In practice, each assumption is recorded as a row
\begin{equation}
    a_i =
    (
        \text{module},
        \text{assumption id},
        \text{severity},
        \text{blocking-for-3D},
        \text{current assumption},
        \text{3D requirement}
    ),
\end{equation}
which is exactly the structure emitted by
\texttt{\_scalar\_assumption\_rows()}.

The fake vector-valued corpus in \texttt{dimensional\_lift\_audit.py} is not a
full 3D physics benchmark. It is a contract smoke test. The code constructs
vector measurements \(z_t \in \mathbb{R}^3\), computes vector-compatible summary
features such as path length and displacement norm,
\begin{align}
    \text{path length} &= \sum_{t=2}^{T} \lVert z_t - z_{t-1} \rVert_2, \\
    \text{displacement norm} &= \lVert z_T - z_1 \rVert_2,
\end{align}
and then verifies that the shared prediction/posterior artifact schema still
holds.

\section{Assumptions Behind The Gate Logic}

The advanced-filter decision layer assumes:
\begin{itemize}[leftmargin=1.5em]
    \item simpler methods should be exhausted before more complex ones are added,
    \item switching evidence should be measured rather than asserted,
    \item sensing-limited failures should not be mistaken for inference-limited failures,
    \item RBPF requires a genuine sampled latent structure plus a conditionally
    tractable continuous substate.
\end{itemize}

\section{Advanced Filter Decision Gates}

The main decision surfaces are:
\begin{itemize}[leftmargin=1.5em]
    \item \texttt{advanced\_filter\_decision.py}
    \item \texttt{rl\_backend\_decision.py}
    \item advanced-gate portions of \texttt{generic\_filtering\_contract.py}
\end{itemize}

These modules ask what failure evidence would justify:
\begin{itemize}[leftmargin=1.5em]
    \item IMM,
    \item PF,
    \item RBPF.
\end{itemize}

\section{Why IMM, PF, and RBPF Remain Decision-Gated}

The current rule is not merely verbal. The implemented gate logic is based on
measured gains from simpler rungs. For IMM-like switching logic, the current
decision report computes quantities such as
\begin{align}
    \Delta_{\text{post-switch}}^{\text{TM-static}}
    &= A_{\text{post-switch}}^{\text{transition matrix}}
    - A_{\text{post-switch}}^{\text{static accumulator}}, \\
    \Delta_{\text{post-switch}}^{\text{TM-kalman}}
    &= A_{\text{post-switch}}^{\text{transition matrix}}
    - A_{\text{post-switch}}^{\text{switching Kalman bank}}.
\end{align}
As long as the transition-matrix accumulator still buys positive switching
improvement and the switching Kalman bank has not clearly displaced it, the repo
does not yet have evidence that IMM complexity is necessary.

For particle filtering, the current decision surface is driven by the absence of
a demonstrated nonlinear or non-Gaussian benchmark and by the fact that the
strongest hard case remains sensing-limited. The implemented evidence rows in
\texttt{advanced\_filter\_decision.py} are therefore evaluating a logic of the
form
\begin{align}
    \text{PF justified}
    &\Longrightarrow
    \text{nonlinear benchmark exists}
    \land
    \text{robust Kalman still fails}
    \notag \\
    &\qquad \land
    \text{failure is not primarily sensing-limited}.
\end{align}
The artifact
\texttt{artifacts/advanced\_filter\_decision\_v1/advanced\_filter\_decision\_numeric\_walkthrough.md}
now expands one real decision pass numerically. It substitutes the measured
transition gains, identifiability gaps, velocity-aided improvement, and best
Kalman outlier accuracy into the gate interpretation so the current
``defer IMM / defer PF'' result is justified by actual evidence values rather
than only a prose rule.

For RBPF, the required latent structure is stricter still. The intended
decision criterion is:
\begin{equation}
    \text{RBPF justified}
    \Longrightarrow
    x_t = (u_t, v_t),
\end{equation}
with \(u_t\) representing a sampled nonlinear or discrete latent structure and
\(v_t\) representing a conditionally linear-Gaussian substate that can still be
filtered analytically.

\section{Interpretation Of The Current Gate}

The current rule is therefore:
\begin{itemize}[leftmargin=1.5em]
    \item IMM is deferred until switching evidence clearly exceeds what the
    current transition-aware ladder handles.
    \item PF is deferred until nonlinear or non-Gaussian behavior defeats the
    current recursive methods for reasons other than sensing limits.
    \item RBPF is deferred until there is a real sampled-latent plus
    conditionally tractable continuous-state split.
\end{itemize}

That is why the advanced-filter decision reports are architecture artifacts, not
mere TODO notes. They encode the proof obligation for adding more complexity.

\section{Artifact Map}

Primary artifacts for this layer:
\begin{itemize}[leftmargin=1.5em]
    \item \texttt{artifacts/dimensional\_lift\_audit/}
    \item \texttt{artifacts/advanced\_filter\_decision\_v1/}
    \item \texttt{artifacts/filtering\_contract/}
\end{itemize}

\end{document}
